% Created 2026-08-06 gio 09:21
% Intended LaTeX compiler: pdflatex
\documentclass[11pt]{article}
\usepackage[utf8]{inputenc}
\usepackage[T1]{fontenc}
\usepackage{graphicx}
\usepackage{longtable}
\usepackage{wrapfig}
\usepackage{rotating}
\usepackage[normalem]{ulem}
\usepackage{amsmath}
\usepackage{amssymb}
\usepackage{capt-of}
\usepackage{hyperref}
\usepackage[newfloat]{minted}
\author{marco robba}
\date{\today}
\title{}
\hypersetup{
 pdfauthor={marco robba},
 pdftitle={},
 pdfkeywords={},
 pdfsubject={},
 pdfcreator={Emacs 30.2 (Org mode 9.7.29)}, 
 pdflang={English}}
\begin{document}

\tableofcontents

\section{Ripasso UML differenza tra diagramma a stati e di sequenza}
\label{sec:orgc1df1d3}
nota : diobastardo comprendere in modo completo tutti questi livelli è infinito
\subsection{Scenari Vs RF}
\label{sec:org7ff3d43}
Un RF è una dichiarazione generalistica di cosa il sistema deve fare, è astratta, lo scenario è un caso particolare che può effettuare il sistema
Uno sviluppo gerarchico degli RF porta a la trattazione di tutti casi possibili secondo condizioni fino a quelle atomiche, definire dunque in questa maniera attraverso tale albero tutte i casi d'uso e le condizioni e le situazioni che si devono gestire per implementarli
una dipendenza è una specificazione, la gestione dei guasti dipende dalla gestione se un treno si trova dentro una tratta o no e a sua volta si specializza nella gestione se un treno si trova su una tratta o no
Il criterio è il test di eliminazione: A dipende da B se togliendo B il requisito A non è più soddisfacibile. Prova:
Gli scenari servono:
\begin{enumerate}
\item ad aumentare la comprensione del sistema per la modellazione dei diagrammi
\item far comprendere immediatamente al lettore cosa è capace di fare il letto
\end{enumerate}
Attenzione: il titolo di uno scenario può tranquillamente essere un RF

In sintesi, i requisiti funzionali sono le regole del gioco, mentre gli scenari sono le partite giocate per dimostrare che quelle regole funzionano.


\begin{itemize}
\item Scenari ancora più in particolare:
\end{itemize}
La regola d'oro è: si modellano gli scenari significativi e rappresentativi, che coprono le casistiche più importanti. In pratica, si segue questa strategia:

Lo Scenaro Principale (``Happy Path'' / Percorso Felice): Quello che funziona tutto bene, senza intoppi. È la sequenza ``standard'' che l'utente medio segue per raggiungere il suo obiettivo.

Gli Scenari Alternativi (``Sad Path'' / Percorsi Tristi): Le varianti più comuni o più critiche. Ad esempio: ``Cosa succede se la password è sbagliata?'', ``Cosa succede se l'utente clicca su `Recupera password'?''.

Gli Scenari di Eccezione (``Corner Cases'' / Casi limite): Quelli che mettono alla prova i confini del sistema. Ad esempio: ``Cosa succede se il campo email è vuoto?'', ``Cosa succede se l'utente ha il profilo bloccato?''.

Gli Scenari ad Alto Rischio: Quelli che, se falliscono, causerebbero danni enormi (perdita di dati, problemi di sicurezza, violazioni legali). Questi vanno sempre modellati, anche se sono rari.

La strategia vincente: ``Copertura, non esaustività''
\subsection{UML differenza tra diagramma a stati e di sequenza}
\label{sec:org9c081db}

nel diagramma i messaggi tra i vari componenti si possono mettere ma sarebbero dei normali cambi di stato e non vengono messi in risalto
nel diagramma di secuenza gli stati di un programma: i puti di esecuzione del codice divisi per il  salto stesso che possono fare non sono messi in risalto perche non mostra i salti fra gli stati
inoltre fare un albero tipo:
A
\begin{itemize}
\item B
\item C
\end{itemize}
idica che A si specifica / dipende dalle seguenti distinte parti dunque tutto l'edge dipende da D quindi anche se tutte le parti di A dipendono da D allora A dipende da D A-->D
la matematica delle dipendenze è delle generalizzazioni è tosta.
Cosa dice Gemini
\begin{enumerate}
\item Diagramma di Sequenza (Sequence Diagram)

Domanda a cui risponde: ``Chi parla con chi, e in quale ordine temporale?''

Cosa mette in risalto: Le linee di vita (lifelines) degli oggetti, la sequenza dei messaggi (chiamate di metodi, eventi, risposte) e la durata delle attivazioni.

Perché ignora il codice interno: Se dovessi mostrare ogni if, switch, salto di codice o cambio di stato interno, il diagramma diventerebbe illeggibile. Per gestire i salti condizionali si usano al massimo i frammenti combinati (come alt, opt, loop), ma sempre mantenendo l'attenzione sui messaggi scambiati tra oggetti.

\item Diagramma di Stato (Statechart Diagram)

Domanda a cui risponde: ``In che stato si trova il sistema o l'oggetto e come reagisce agli eventi?''

Cosa mette in risalto: Gli stati (es. In attesa, In elaborazione, Errore) e le transizioni causate da eventi o condizioni.

Perché i messaggi sono secondari: I messaggi o le azioni intermedie esistono solo come ``effetti collaterali'' delle transizioni, ma il protagonista assoluto è l'evoluzione dello stato dell'entità.
\end{enumerate}
\subsection{Esercizio}
\label{sec:orgc9da7a3}
\begin{center}
\includegraphics[width=.9\linewidth]{/tmp/babel-m8y22I/plantuml-oJ2Swk.png}
\label{}
\end{center}
\begin{center}
\includegraphics[width=.9\linewidth]{/tmp/babel-m8y22I/plantuml-Ve1zx7.png}
\label{}
\end{center}
\section{nuovo (è incompleto il resto è nella doc finale)}
\label{sec:org2604714}
Basato da dove parte l'evento
\subsection{RF}
\label{sec:org7189aed}
Il numero organico è troppo alto provare una riduzione col grafo delle dipendenze
aggiungendo le dipendenze si possono aggiungere le specificazioni a diamante
RF01: casi d'uso (dipende da  Correttezza del sistema gestionale)
\begin{itemize}
\item visualizzazione dei dati in tempo reale e corretti
\item visualizzazione dei guasti dipende
\item modifica dei dati da parte esclusiva dell'aministratore
\item costruzione di uno storico (funzionalità non richiesta)
\end{itemize}
RF02: Correttezza del sistema gestionale
\begin{itemize}
\item gestione guasti(stazioni, treni)
\begin{itemize}
\item guasto stazione
\begin{itemize}
\item guasto di perdità di connessione
\begin{itemize}
\item chashare i dati prodotti fino alla ripresa della nuova connessione quali verranno reinviati in massa
\end{itemize}
\item guasto che ne impedisce la percorrenza dei treni
\begin{itemize}
\item sensibile al treno il treno si deve fermare se la prossima stazione ha il guasto
\end{itemize}
\end{itemize}
\item guasto treno
\begin{itemize}
\item memorizzarlo nel db
\item gestione alle conseguenze
\begin{itemize}
\item g. sensibilità alla stazione stazione: il treno si guasta nella stazione si vuole che la stazione risulterà in'usabile(Guasta)
\item g. sensibile tratta: se il treno si guasta nella tratta si vuole che la stazione trattà risulterà in'usabile(Guasta)
\end{itemize}
\end{itemize}
\end{itemize}
\item gestione treni
\begin{itemize}
\item gestione ritardi
\begin{itemize}
\item se il treno è fermo quando dovrebbe circolare accumulo del ritardo
\begin{itemize}
\item per guasto (dipende da gestione dei guasti)
\item per ordinamento dell'amiministratore.(non implementato)
\end{itemize}
\end{itemize}
\item servire le informazioni di stato (simulazione del treno)(dipende dalla gestione dei guasti)
\begin{itemize}
\item posizione puramente a quale tratta si trova/stazione
\item ritardo (dipende da 'gestione ritardi')
\end{itemize}
\end{itemize}
\item gestione stazioni
\begin{itemize}
\item informazioni da rendere alla visualizzazione
\begin{itemize}
\item attraverso la gestione della chash
\item treni attualmente presenti nella stazione
\item treno e tempo di entrata nella stazione
\item treno e tempo di uscita dalla stazione
\end{itemize}
\end{itemize}
\item Validazione nodo: si vuole che il nodo una volta entrato nel sistema il proprio id che dichiara venga
validato attraverso la voce del db
\begin{itemize}
\item se è presente è ok, se non è presente non può entrare nel sistema
\end{itemize}
\item modifica tratta
\begin{itemize}
\item il treno non è fuori dalla tratta il treno viene fermato in attesa di essere soppresso
\item il treno è dentro la tratta, la tratta deve essere seguita dal punto in cui è il treno
\end{itemize}
\end{itemize}
\section{precedente}
\label{sec:org3a61191}
\subsection{microservizio treno}
\label{sec:orgb9f858f}
\begin{itemize}
\item aggiornamento tratta
\begin{itemize}
\item salvarla in ram per il proprio uso e sul db
\end{itemize}
\item gestione causa guasti alla prossima fermata stazione o binario (tratatta ovvero copia di stazioeni)
\begin{itemize}
\item ricezione dai sesori --> modifica della percorrenza della tratta, sincronizzazione sui topic di allarme delle
nuove stazioni
\begin{itemize}
\item se arrivano evendi del tipo la prossima stazione è fuori uso aggiornare variabile in rame e il database
che il treno si è femato e
\item se l'evento è il ripristino della funzionalità della stazione allora riaggiornare la ram e il db che il
treno è partito
\end{itemize}
\end{itemize}
\item Ritardi la tratta possiede orari corretti di percorrenza
\item inviare in certrale l'entrata e uscita da una stazione topic treni/+/passaggio
formato : boh lo fa l'ai
\item invio se il treno si guasta alla centrale
\end{itemize}
\subsection{microservizio stazione}
\label{sec:org08bfbdd}
\begin{itemize}
\item invio alla centrale i dati del passaggio dei treni
\begin{itemize}
\item i sensori della stazione informano id treno e se è in entrata o in uscita,
poi si aggiorna la ram e si inviano i dati alla stazione centrale per la storicizzazione
\end{itemize}
\item la staione può perdere la connessione
 con la centrale nel caso seguire il diagramma degli stati che ne
descrive la dinamica di gestione
\end{itemize}
\subsection{Centrale}
\label{sec:org6d76551}
\begin{itemize}
\item deve rispondere alle richieste di visualizzazione
\begin{itemize}
\item montare su il motore in tepo reale, cicli continui che inviano i dati se la richiesta è ancora attiva
\end{itemize}
\item formattare e storicizzare i dati
\end{itemize}
L'unica crud a cui serve la propagazinone ai microservizzi è la
\begin{itemize}
\item gestione guasti
\begin{itemize}
\item se i microservizzi avvertono guasti , aggiornamento del db,
\item se il tecnico invia l'operatore aggiornare il db di modo che ormai il guasto sia terminato e
pubblicare sul topic degli allert generale per tutti i treni o tutte le stazioni che quel tale microservizio è stato risolto
\end{itemize}
\end{itemize}
interrogare il database ogni 10 secondi per ogni tipo di visualizzazione di modo da rendere il ``tempo reale''
\end{document}
